% ---------- Datos del documento ----------
\newcommand{\tituloSA}{Medir lo que no se puede alcanzar: ¿dónde caben más paneles solares en nuestro instituto?}
\newcommand{\autor}{Javier Gómez López}
\newcommand{\asignatura}{Diseño Curricular de las Matemáticas}
\newcommand{\master}{Máster Universitario en Formación del Profesorado -- Especialidad en Matemáticas}

% ---------- Paquetes ----------
\usepackage{float}
\usepackage[utf8]{inputenc}
\usepackage[T1]{fontenc}
\usepackage[spanish,es-noshorthands,es-tabla]{babel}
\usepackage[scaled=0.92]{helvet} % equivalente libre de la Arial del original
\renewcommand{\familydefault}{\sfdefault}
\usepackage[left=2cm,right=2cm,top=2.55cm,bottom=1.6cm,
            headheight=44pt,headsep=4pt,footskip=20pt]{geometry}
\usepackage{amsmath,amssymb}
\usepackage{xcolor}
\usepackage{graphicx}
\usepackage{fancyhdr}
\usepackage{lastpage}
\usepackage{etoolbox}
\usepackage{enumitem}
\usepackage{titlesec}
\usepackage{pdflscape}
\usepackage{tikz}
\usetikzlibrary{angles,quotes,arrows.meta}
\usepackage{tabularray}
\usepackage[most]{tcolorbox}
\usepackage{csquotes}
\usepackage[backend=bibtex,style=authoryear,maxcitenames=2,uniquename=false,uniquelist=false]{biblatex}
\addbibresource{bibliografia.bib}
\usepackage[hidelinks]{hyperref}

\setlength{\parindent}{0pt}
\setlength{\parskip}{5pt}
\setlist{nosep,leftmargin=1.2em}
\emergencystretch=1.5em

% ---------- Colores (medidos sobre la plantilla) ----------
\definecolor{urjc}{HTML}{C8102E}
\definecolor{rosa}{HTML}{FBE9EB}
\definecolor{borde}{HTML}{D0D0D0}
\definecolor{titular}{HTML}{3C3C3B}
\definecolor{pistagris}{HTML}{7F7F7F}

% ---------- Títulos ----------
\titleformat{\section}{\Large\bfseries\color{urjc}}{\thesection}{0.6em}{}[{\color{urjc}\titlerule[0.7pt]}]
\titleformat{\subsection}{\large\bfseries\color{titular}}{\thesubsection}{0.6em}{}
\titleformat{\subsubsection}{\normalsize\bfseries\color{urjc}}{}{0pt}{}
\titlespacing*{\section}{0pt}{14pt}{6pt}
\titlespacing*{\subsection}{0pt}{10pt}{3pt}
\titlespacing*{\subsubsection}{0pt}{8pt}{2pt}
\setcounter{secnumdepth}{2}
\setcounter{tocdepth}{2}
\makeatletter\renewcommand*\l@subsection{\@dottedtocline{2}{1.5em}{2.9em}}\makeatother

% ---------- Cabecera y pie ----------
\newcommand{\cabecera}{Máster en Formación del Profesorado · Diseño Curricular de las Matemáticas · 2026-27}

\pagestyle{fancy}
\fancyhf{}
\fancyhead[L]{\includegraphics[height=42pt]{logo_urjc.png}}
\fancyhead[R]{\footnotesize\color{pistagris}\cabecera}
\fancyfoot[R]{\footnotesize\color{pistagris}Página \thepage\ de \pageref{LastPage}}
\renewcommand{\headrule}{{\color{urjc}\hrule height 1pt}}
\renewcommand{\footrulewidth}{0pt}

% ---------- Elementos de texto ----------
% Etiqueta de un campo (Título:, Producto:...)
\newcommand{\campo}[1]{{\bfseries\color{urjc}#1}}
% Texto orientativo, en cursiva gris
\newcommand{\pista}[1]{{\itshape\small\color{pistagris}#1}}
% Lista con viñetas; dentro de las tablas hay que usar esta y no itemize
\newenvironment{lista}
  {\begin{minipage}[t]{\linewidth}\begin{itemize}[leftmargin=*,nosep]}
  {\end{itemize}\end{minipage}}
% Encabezado de los apartados a todo el ancho de la plantilla (7 y 8)
\newcommand{\apartado}[1]{%
  \par\vspace{14pt}%
  {\large\bfseries\color{urjc}#1\par}%
  \vspace{4pt}{\color{urjc}\hrule height 0.7pt}\vspace{5pt}}

% ---------- Recuadros ----------
% \begin{tarjeta}[subtítulo opcional]{título}
\newtcolorbox{tarjeta}[2][]{%
  enhanced, sharp corners, boxrule=0.5pt, titlerule=0pt,
  colframe=borde, colback=white, colbacktitle=urjc, coltitle=white,
  fonttitle=\bfseries\large, toptitle=3pt, bottomtitle=3pt,
  left=7pt, right=7pt, top=7pt, bottom=7pt,
  halign=flush left, halign title=flush left,
  before upper={\setlength{\parskip}{5pt}},
  title={#2\ifstrempty{#1}{}{\\[-1pt]{\footnotesize\itshape #1}}}}

% Fila de dos recuadros de igual altura
\newenvironment{filatarjetas}{%
  \begin{tcbraster}[raster columns=2, raster equal height=rows,
                    raster column skip=8pt, raster after skip=12pt]}
  {\end{tcbraster}}

% Sesión de la secuencia: \sesion{n}{tipo}{título}{duraciones y texto}
\newtcolorbox{sesionbox}[3]{%
  enhanced, sharp corners, boxrule=0.5pt, titlerule=0pt,
  colframe=borde, colback=white, colbacktitle=urjc, coltitle=white,
  fonttitle=\bfseries, toptitle=2pt, bottomtitle=2pt,
  left=6pt, right=6pt, top=4pt, bottom=5pt, after skip=8pt,
  halign=flush left, halign title=flush left,
  fontupper=\small, before upper={\setlength{\parskip}{3pt}},
  title={Sesión #1 · #2: #3}}

% Caja de información
\newtcolorbox{nota}[1]{%
  enhanced, sharp corners, boxrule=0pt, leftrule=3pt, colframe=urjc, colback=rosa,
  coltitle=urjc, colbacktitle=rosa, fonttitle=\bfseries, attach title to upper={\ },
  left=7pt, right=7pt, top=4pt, bottom=4pt, title={#1}, fontupper=\small}

% ---------- Tablas ----------
\NewTblrEnviron{tablaSA}
\SetTblrInner[tablaSA]{
  hlines={0.5pt,borde}, vlines={0.5pt,borde},
  leftsep=5pt, rightsep=5pt, rowsep=3pt,
  cells={font=\footnotesize, valign=t},
  row{1}={bg=rosa, fg=urjc, font=\footnotesize\bfseries},
}

% Versión con salto de página y cabecera repetida
\NewTblrTheme{limpio}{
  \SetTblrTemplate{caption}{empty}
  \SetTblrTemplate{capcont}{empty}
  \SetTblrTemplate{contfoot}{empty}
  \SetTblrTemplate{conthead}{empty}
  \SetTblrTemplate{note}{empty}
  \SetTblrTemplate{remark}{empty}
}
\NewTblrEnviron{tablaSAlarga}
\SetTblrOuter[tablaSAlarga]{long, theme=limpio, entry=none, label=none, presep=2pt, postsep=4pt}
\SetTblrInner[tablaSAlarga]{
  hlines={0.5pt,borde}, vlines={0.5pt,borde},
  leftsep=5pt, rightsep=5pt, rowsep=3pt,
  cells={font=\footnotesize, valign=t},
  rowhead=1,
}
